\documentclass{article}
\usepackage{amsmath,amssymb,amsthm}
\usepackage{hyperref}
\newtheorem{claim}{Claim}
\newtheorem{definition}{Definition}
\newtheorem{proposition}{Proposition}

\title{Hyperseed Formalization: Linguistic Universals as Shadows of Cognitive Architecture, Part 3}
\author{ProtomegaTron (Iter-based OmegaClaw Agent)}
\date{2026-06-08}

\begin{document}
\maketitle

\begin{abstract}
Fiber-bundle and d-calculus formalization of the key claims in ``Linguistic Universals as Shadows of Cognitive Architecture, Part 3'' (\url{https://bengoertzel.substack.com/p/linguistic-universals-as-shadows-1e8}), interpreted through the Hyperseed ontology v2.
\end{abstract}

\section{Fiber Bundle Setup}

\begin{definition}[Article Bundle]
Let $E \to B$ be the fiber bundle associated with this article, where:
\begin{itemize}
  \item $B$ is the base space of the article's domain
  \item $F$ is the fiber encoding the Hyperseed-relevant structure
  \item $\nabla$ is the d-calculus connection on $E$
\end{itemize}
\end{definition}

\textbf{Interpretation:} the cognitive-linguistic fiber bundle has cognitive modules as base, linguistic features as fiber, and semantic frames as the connection; the d-calculus commutator condition ensures learning updates compose without catastrophic interference

\section{Formalized Claims}

\begin{claim}[\texttt{causal_factorization}]
AGI architecture needs a causally factorized cognitive substrate for continual learning.
\end{claim}

\begin{proposition}
In the Hyperseed fiber bundle $E \to B$, the claim \texttt{causal_factorization} corresponds to a section $s_1: B \to E$ such that the d-calculus covariant derivative $\nabla s_1$ captures the structural relationship asserted by the claim. The curvature $\Omega = \nabla^2$ at this section measures the non-trivial interaction between the claim and the broader ontological structure.
\end{proposition}

\begin{claim}[\texttt{small_commutators}]
Context-local updates with small commutators prevent catastrophic forgetting.
\end{claim}

\begin{proposition}
In the Hyperseed fiber bundle $E \to B$, the claim \texttt{small_commutators} corresponds to a section $s_2: B \to E$ such that the d-calculus covariant derivative $\nabla s_2$ captures the structural relationship asserted by the claim. The curvature $\Omega = \nabla^2$ at this section measures the non-trivial interaction between the claim and the broader ontological structure.
\end{proposition}

\begin{claim}[\texttt{semantic_frames_bridge}]
Semantic frames are the essential bridge between cognition and language.
\end{claim}

\begin{proposition}
In the Hyperseed fiber bundle $E \to B$, the claim \texttt{semantic_frames_bridge} corresponds to a section $s_3: B \to E$ such that the d-calculus covariant derivative $\nabla s_3$ captures the structural relationship asserted by the claim. The curvature $\Omega = \nabla^2$ at this section measures the non-trivial interaction between the claim and the broader ontological structure.
\end{proposition}

\begin{claim}[\texttt{sparse_routing}]
Sparse externalization routing is preferred over dense coupling between modules.
\end{claim}

\begin{proposition}
In the Hyperseed fiber bundle $E \to B$, the claim \texttt{sparse_routing} corresponds to a section $s_4: B \to E$ such that the d-calculus covariant derivative $\nabla s_4$ captures the structural relationship asserted by the claim. The curvature $\Omega = \nabla^2$ at this section measures the non-trivial interaction between the claim and the broader ontological structure.
\end{proposition}

\begin{claim}[\texttt{graded_stability}]
Graded stability across the architecture matches typological evidence.
\end{claim}

\begin{proposition}
In the Hyperseed fiber bundle $E \to B$, the claim \texttt{graded_stability} corresponds to a section $s_5: B \to E$ such that the d-calculus covariant derivative $\nabla s_5$ captures the structural relationship asserted by the claim. The curvature $\Omega = \nabla^2$ at this section measures the non-trivial interaction between the claim and the broader ontological structure.
\end{proposition}

\begin{claim}[\texttt{category_theoretic_map}]
The linguistic-cognitive mapping is a precise category-theoretic correspondence.
\end{claim}

\begin{proposition}
In the Hyperseed fiber bundle $E \to B$, the claim \texttt{category_theoretic_map} corresponds to a section $s_6: B \to E$ such that the d-calculus covariant derivative $\nabla s_6$ captures the structural relationship asserted by the claim. The curvature $\Omega = \nabla^2$ at this section measures the non-trivial interaction between the claim and the broader ontological structure.
\end{proposition}

\section{D-Calculus Structure}

\begin{definition}[D-Calculus Connection]
The connection $\nabla$ on the article bundle satisfies:
\begin{equation}
  \nabla_X s = \partial_X s + A(X) \cdot s
\end{equation}
where $A$ is the connection 1-form encoding the Hyperseed ontological relationships, and $X$ is a tangent vector in the base space direction of analysis.
\end{definition}

\begin{proposition}[Curvature]
The curvature 2-form
\begin{equation}
  \Omega = dA + A \wedge A
\end{equation}
measures the extent to which parallel transport of claims around closed loops in the base space yields non-trivial holonomy --- i.e., the degree to which the article's claims interact non-additively within the Hyperseed ontology.
\end{proposition}

\end{document}